%% ============================================================
%%  SECCIÓN 6 — Procedimiento Experimental
%%  Responsable: Minilux
%% ============================================================

\section{Procedimiento Experimental}

El procedimiento se dividió en dos etapas independientes, correspondientes
a cada método de determinación del coeficiente de tensión superficial.

\subsection{Primer método — Balanza de Mohr--Westphal con anillo}

\begin{enumerate}
    \item Se midieron con el vernier los radios externo e interno del anillo
    metálico, registrando $R_{\text{ext}} = \qty{5.040 \pm 0.005}{\centi\meter}$ y
    $R_{\text{int}} = \qty{4.980 \pm 0.005}{\centi\meter}$.

    \item Se suspendió el anillo del brazo mayor de la balanza de
    Mohr--Westphal y se colocó el vasito de plástico vacío en el extremo
    del brazo menor. Se ajustó el contrapeso hasta obtener el equilibrio
    con el brazo mayor en posición horizontal.

    \item Se sumergió el anillo en el recipiente con agua, asegurando que
    quedara paralelo a la superficie líquida y formara una película
    continua a lo largo de todo su perímetro.

    \item Se añadió arena muy lentamente con la pipeta al vasito de
    plástico hasta que la película se desprendió y el anillo se elevó.
    En ese instante se detuvo la adición de arena.

    \item Se retiró el recipiente con agua sin perturbar el sistema.
    A continuación se colocaron jinetillos en las ranuras numeradas del
    brazo mayor hasta recuperar el equilibrio con el contrapeso,
    registrando la masa de cada jinetillo y su posición (número de
    división).

    \item Se repitieron los pasos 2 a 5 cuatro veces más, variando la
    combinación de jinetillos en cada repetición, hasta completar un total
    de 5 mediciones independientes.
\end{enumerate}

\subsection{Segundo método — Película jabonosa con hilo}

\begin{enumerate}
    \item Se preparó una solución jabonosa con agua y una pequeña cantidad
    de detergente en un recipiente limpio.

    \item Se sumergió completamente el dispositivo de tubitos con hilo en
    la solución jabonosa y se retiró lentamente, asegurando la formación
    de una película continua entre los dos tubos.

    \item Se suspendió el sistema por el alambre del tubo superior,
    manteniéndolo en posición horizontal, y se esperó a que la película
    alcanzara el equilibrio estático.

    \item Con el vernier se midieron las dimensiones geométricas de la
    película: el semiancho del tubo $a$, el semiancho mínimo de la
    curvatura del hilo $b$ y la semialtura de la película $h$.

    \item Se determinó el peso del tubo inferior mediante la balanza de
    precisión, obteniendo su masa $m$ y calculando $P = mg$.

    \item Con los valores de $a$, $b$, $h$ y $P$ se calculó el coeficiente
    de tensión superficial aplicando la Ec.~\eqref{eq:gamma_m2}.
\end{enumerate}